%! Author = sriadityavedantam
%! Date = 3/25/26

\section{Properties of Infinite Series}

\lesson{13}{Friday 12 September 2025 9:10}{Day 13}

\begin{definition}[Infinite Series]
    Given a sequence $(a_1,a_2,a_3,\ldots)$, consider:
    \[
        \sum_{n=1}^\infty a_n.
    \]
\end{definition}

\begin{definition}[Convergence of Series]
    The series $\sum_{n=1}^\infty a_n$ converges if the sequence of partial sums
    \[
        S_N \coloneqq \sum_{n=1}^N a_n
    \]
    converges.
    If $S_N\to L$ as $N\to \infty$, then $\sum_{n=1}^\infty a_n = L$.
\end{definition}

\begin{definition}[Cauchy Criterion]
    The series $\sum a_n$ converges if for all $\eps > 0$, there exists $N\in\n$ such that for all $m\geq n\geq N$,
    \[
        \abs{\sum_{k=n+1}^m a_k} < \eps.
    \]
\end{definition}

\begin{theorem*}{
    If $\sum a_n$ converges, then $a_n\to 0$.
}
\end{theorem*}

\begin{definition}[Harmonic Series]
    \[
        \sum_{n=1}^\infty \dfrac{1}{n}
    \]
    diverges.
\end{definition}

\begin{theorem*}{
    If all $a_n\geq 0$, then $\sum a_n$ converges if and only if the sequence of partial sums $(S_N)$ is bounded.
}
\end{theorem*}

\begin{definition}[Geometric Series]
    A series is geometric if for some $r\neq 0$ and $a_0\neq 0$,
    \[
        \dfrac{a_{n+1}}{a_n} = r.
    \]
\end{definition}

\begin{remark}
    That is
    \[
        a_0 + ra_0 + r^2 a_0 + \ldots = a_0(1 + r + r^2 + \ldots).
    \]
    Also take $r = \dfrac{1}{2}$:
    \[
        1 + \dfrac{1}{2} + \dfrac{1}{4} + \dfrac{1}{8} + \ldots = 2.
    \]
\end{remark}

\begin{lemma}{
    If $\abs{r} < 1$, then $r^n\to 0$.
}
    This is equivalent to $\abs{r}^n\to 0$.
    For $0 < r < 1$, we have
    \[
        1 = r^0 > r > r^2 > r^3 > \ldots \geq 0.
    \]
    Thus $(r^n)$ is bounded and monotonic, so it converges.
    Let $r^n \to L$.
    Then $r^{n+1}\to L$ implies $rL = L$, so $(1-r)L = 0$.
    Since $r\neq 1$, we have $L = 0$.
\end{lemma}

\begin{lemma}{
    \[
        \sum_{n=0}^N r^n = \dfrac{1 - r^{N+1}}{1-r}.
    \]
}
    \begin{align*}
        \left(\sum_{n=0}^N r^n\right) (1-r)
        &= (1 + r + r^2 + \ldots + r^N)(1 - r)\\
        &= (1 + r + r^2 + \ldots + r^N) - (r + r^2 + \ldots + r^N + r^{N+1})\\
        &= 1 - r^{N+1}.
    \end{align*}
    Hence
    \[
        \sum_{n=0}^N r^n = \dfrac{1 - r^{N+1}}{1-r}.
    \]
\end{lemma}

\begin{theorem}[Geometric Series Convergence]{
    $\sum r^n$ converges if and only if $\abs{r} < 1$.
}
    If $\abs{r}\geq 1$, then $r^n \not\to 0$, so the series diverges.
    If $\abs{r} < 1$, then
    \begin{align*}
        S_N = \dfrac{1 - r^{N+1}}{1-r}
        &\to \dfrac{1}{1-r} - \dfrac{\lim\limits_{N\to\infty} r^{N+1}}{1 -r}\\
        &= \dfrac{1}{1-r}.
    \end{align*}
    Hence the series converges.
\end{theorem}

\begin{definition}[Absolute Convergence]
    If $\sum \abs{a_n}$ converges, then $\sum a_n$ converges absolutely.
\end{definition}

\begin{definition}
    If $\sum a_n$ converges but $\sum \abs{a_n}$ diverges, then the series converges conditionally.
\end{definition}

\begin{theorem}{
    If $\sum a_n$ converges absolutely, then it converges.
}
    Apply the Cauchy Criterion.
    Let $\eps > 0$.
    Then there exists $N\in\n$ such that for $m > n > N$,
    \[
        \sum_{k=n+1}^m \abs{a_k} < \eps.
    \]
    Thus
    \[
        \abs{\sum_{k=n+1}^m a_k}\leq \sum_{k=n+1}^m \abs{a_k} < \eps.
    \]
\end{theorem}

\begin{example}
    \begin{gather*}
        \sum_{n=1}^\infty \dfrac{(-1)^n}{n} \quad\text{converges conditionally}\\
        \sum_{n=1}^\infty \dfrac{(-1)^n}{n^2}\quad\text{converges absolutely}\\
    \end{gather*}
    Any convergent geometric series converges absolutely.
    \[
        \sum \abs{r^n} = \sum \abs{r}^n = \dfrac{1}{1 - \abs{r}}.
    \]
\end{example}